\documentclass[11pt,a4paper]{article}

\usepackage[margin=1in]{geometry}
\usepackage{amsmath}
\usepackage{amssymb}
\usepackage{amsthm}
\usepackage{algorithm}
\usepackage{algpseudocode}
\usepackage{booktabs}
\usepackage{array}
\usepackage{enumitem}
\usepackage{xcolor}
\usepackage{microtype}
\usepackage[hidelinks]{hyperref}

\DeclareMathOperator*{\argmax}{arg\,max}
\DeclareMathOperator*{\lexmax}{lex\,max}

\newtheorem{proposition}{Proposition}
\newtheorem{theorem}{Theorem}
\newtheorem{lemma}{Lemma}
\theoremstyle{definition}
\newtheorem{definition}{Definition}
\theoremstyle{remark}
\newtheorem{remark}{Remark}

\newcommand{\ind}[1]{\mathbf{1}\!\left[#1\right]}
\newcommand{\canplace}{\textsc{CanPlace}}
\newcommand{\score}{\textsc{Score}}
\newcommand{\nightsolve}{\textsc{NightSolve}}
\newcommand{\weeksolve}{\textsc{WeekSolve}}
\newcommand{\canon}{\textsc{Canon}}
\newcommand{\lex}{\succ_{\mathrm{lex}}}

\algrenewcommand\algorithmicrequire{\textbf{Input:}}
\algrenewcommand\algorithmicensure{\textbf{Output:}}

% Ragged-right paragraph columns: narrow measures justify badly.
\newcolumntype{L}[1]{>{\raggedright\arraybackslash}p{#1}}

\title{\bfseries The Grave-Shift Placement Problem\\[0.35em]
\large A Lexicographic Assignment Formulation}
\author{ShiftBuilder --- Internal Maintenance, Gun Lake Casino}
\date{18 July 2026}

\begin{document}
\maketitle

\begin{abstract}
\noindent
We give a formal specification of the nightly team-member placement problem
solved by ShiftBuilder, and an algorithm for it. The operational requirement
is a strict priority ordering over four objectives---coverage, rotation,
preference, and skill---together with a family of hard constraints that are
never tradeable. We show that this is precisely a \emph{lexicographic
assignment problem}: a maximum-cardinality bipartite matching for the
dominant objective, refined by lexicographic vector comparison over the
remainder. The formulation yields provable coverage optimality, exact
infeasibility certificates via the deficiency form of Hall's theorem, and
structural enforcement of the priority ordering. We contrast it with the
scalarized (weighted-sum) formulation currently deployed, and identify the
classes of defect that the lexicographic formulation renders
unrepresentable.
\end{abstract}

\section{Scope and motivation}

Each night, a roster of approximately twenty-five team members must be
assigned to approximately thirty-four board positions: ten zones, ten
restroom stations, an administrative position, a smoking-room position, a
flexible auxiliary row, and two bands of partial-shift overlap seats. The
assignment is subject to hard physical and policy constraints, and is
evaluated against four objectives whose relative priority has been ratified
by the operator:
\begin{equation}
\label{eq:hierarchy}
\textsf{coverage} \;\succ\; \textsf{rotation} \;\succ\;
\textsf{preference} \;\succ\; \textsf{skill}.
\end{equation}

The symbol $\succ$ in \eqref{eq:hierarchy} denotes \emph{priority}, not a
weight ratio. The operator did not assert that one unit of coverage is worth
some number of units of rotation; the assertion is that coverage decides,
and the remaining objectives are consulted only among coverage-equal
solutions. Section~\ref{sec:objective} makes this precise, and
Remark~\ref{rem:bigm} explains why encoding a priority ordering as a weighted
sum is a lossy translation.

\section{Notation}
\label{sec:notation}

\begin{definition}[Instance]
A placement instance is a tuple
$\mathcal{I} = (T, S, \mathcal{R}, \mathrm{lock}, \mathrm{pin}, H, W, \theta)$
where
\begin{itemize}[itemsep=0.15em,topsep=0.25em]
  \item $T$ is the set of team members scheduled tonight, after removal of
        call-offs and dated unavailability;
  \item $S$ is the ordered set of board positions, indexed by fill order;
  \item $\mathcal{R} \subseteq S$ is the set of \emph{required} positions
        (restrooms and zones); we write $\mathcal{O} = S \setminus
        \mathcal{R}$ for the \emph{optional} remainder;
  \item $\mathrm{lock} \subseteq S$ are operator-locked positions, each with
        an incumbent $\iota(s) \in T$;
  \item $\mathrm{pin} \subseteq S$ are preserved-but-unlocked existing
        placements, per the run's preserve policy;
  \item $H(t)$ is the placement history of $t$, a sequence of
        $(\textit{night}, \textit{area})$ events in decreasing night order;
  \item $W(t)$ is the set of placements already planned for $t$ earlier in
        the current week solve;
  \item $\theta$ is the active engine configuration (weights and thresholds),
        and $\zeta$ a pseudorandom seed.
\end{itemize}
\end{definition}

An \emph{assignment} is a partial injection $\mu : S \rightharpoonup T$;
$\mu(s) = t$ means $t$ holds position $s$. Injectivity encodes the physical
constraint that a person occupies at most one position per night. We write
$\mathrm{dom}(\mu)$ for the filled positions and
$\mathrm{cov}(\mu) = \lvert \mathrm{dom}(\mu) \cap \mathcal{R} \rvert$ for
the coverage of $\mu$.

\subsection{Area identity}

Rotation is not defined over positions but over \emph{areas}. Let
$\alpha : S \to \mathcal{A}$ be the area map. It is not injective: the two
sides of a restroom pair share an area, because working either side of
restroom~8 is operationally the same exposure, and overlap seats within a
band are fungible. Zones and auxiliary roles are their own areas.
\begin{equation}
\alpha(\texttt{MRR8}) = \alpha(\texttt{WRR8}) = \texttt{RR8},
\qquad
\alpha(\texttt{OL-PM-}k) = \texttt{OL-PM},
\qquad
\alpha(\texttt{Z4}) = \texttt{Z4}.
\end{equation}

\subsection{Key normalization}

Position identity is recorded in several vocabularies across the persistence
layer: a user-interface form, a database form, a historical trail form, and
a per-night ephemeral form for flexible auxiliary shells. Let $\canon$ be the
total function mapping any of these to canonical form.

\begin{definition}[Boundary rule]
\label{def:boundary}
Every key entering the algorithm is normalized by $\canon$ at the boundary.
The algorithm is defined only over canonical keys. $\canon$ is total over
production key forms; an unrecognized key yields a distinguished value
$\bot$, and $\canplace(\cdot, \bot, \cdot)$ is false.
\end{definition}

Definition~\ref{def:boundary} is stated as part of the algorithm rather than
as an implementation note because the alternative---comparing keys drawn from
distinct vocabularies---silently produces \emph{false negatives} in every
rotation predicate, and false negatives in a rotation predicate are
indistinguishable from good news. A system in this state reports that a
member has never worked a position they worked last week.

\section{The feasibility gate}
\label{sec:gate}

All hard constraints are collected into a single predicate. No other part of
the algorithm may test eligibility.

\begin{definition}[Gate]
$\canplace(t, s, r)$ holds when member $t$ may occupy position $s$ at
relaxation level $r \in \{0, 1\}$. The constituent gates are:

\smallskip
\noindent\textbf{Level 0} (never relaxed, at any $r$):
\begin{description}[itemsep=0.1em,topsep=0.2em,leftmargin=3.2em,font=\normalfont\itshape]
  \item[$G_1$] Position physics. Zone positions require a full-grave member;
        men's and women's restrooms require the corresponding gender;
        overlap seats require the matching partial-shift band; full-night
        positions exclude partial-shift members. Unknown gender and unknown
        position family both \emph{fail closed}.
  \item[$G_2$] Locks. $s \in \mathrm{lock} \Rightarrow t = \iota(s)$.
  \item[$G_3$] Accommodations, from all recorded sources.
  \item[$G_4$] Schedule membership.
  \item[$G_5$] Operator rules, evaluated \emph{within their declared scope}.
        A rule scoped to a position class constrains only that class.
  \item[$G_6$] Hard avoidance: declared position avoidance, and pairwise
        avoidance with an already-placed neighbour.
\end{description}

\smallskip
\noindent\textbf{Level 1} (relaxed only when $r = 1$):
\begin{description}[itemsep=0.1em,topsep=0.2em,leftmargin=3.2em,font=\normalfont\itshape]
  \item[$G_7$] Rotation. The area $\alpha(s)$ must not occur among the three
        most recent events of $H(t)$ merged with $W(t)$. Continuity roles
        (the administrative position and auxiliary roles flagged as
        continuity) and overlap seats are exempt by policy.
\end{description}
\end{definition}

The ladder has exactly two rungs. There is no level at which a Level-0 gate
yields. Every placement made at $r = 1$ carries a relaxation record in its
provenance, and Section~\ref{sec:ilp} shows how the \emph{number} of such
records is itself minimized.

\begin{remark}
The placement of $G_6$ at Level 0 is a policy choice, and it is the choice
the written operator contract makes: a hard avoidance is described there as a
constraint, not a cost. Should the operator wish avoidance to be tradeable
against coverage, the correct edit is to demote $G_6$ to a term in $\pi$
(Section~\ref{sec:objective}), not to introduce a third relaxation rung.
\end{remark}

\section{The objective}
\label{sec:objective}

\begin{definition}[Score]
For a candidate placement of $t$ at $s$ under partial assignment $\mu$,
\begin{equation}
\score(t, s, \mu) \;=\; \bigl(\rho(t,s),\; \pi(t,s,\mu),\; \kappa(t,s)\bigr)
\in \mathbb{Z} \times \mathbb{R} \times \mathbb{R},
\end{equation}
with components defined as follows.
\end{definition}

\paragraph{Rotation.} Let $\sigma_{30}(t,a)$ be the number of the last thirty
worked nights in which $t$ occupied area $a$, let $\gamma(t,a)$ be the number
of days since $t$ last occupied $a$, and let $\mathrm{sf}$ and $\mathrm{l}_5$
be indicators for a same-side restroom-family repeat within the critical
window and a same-area repeat within the last five events. Then
\begin{equation}
\label{eq:rot}
\rho(t,s) \;=\; 100
\;-\; 55\,\ind{\mathrm{sf}}
\;-\; 25 \min\bigl(\sigma_{30}(t, \alpha(s)),\, 3\bigr)
\;-\; 10\,\ind{\mathrm{l}_5}
\;+\; \min\!\left(7, \left\lfloor \gamma(t,\alpha(s))/8 \right\rfloor\right).
\end{equation}
The prior-three same-area condition does \emph{not} appear in
\eqref{eq:rot}; it is the gate $G_7$, and a gate is not a term.

\paragraph{Preference.} With $h, f \in \mathbb{Z}$ the summed hard and soft
preference signals for $(t,s)$ and $a(t,s,\mu)$ the pair-affinity signal
against already-placed neighbours,
\begin{equation}
\pi(t,s,\mu) \;=\; \mathrm{clamp}(h, 0, 1)
\;+\; 0.4\,\mathrm{clamp}(f, -1, 1)
\;+\; 0.5\,a(t,s,\mu).
\end{equation}
Every term is clamped, so $\pi$ is bounded by construction. Hard
\emph{avoidance} does not appear here; it is $G_6$.

\paragraph{Skill.} With $\varsigma(t)$ the skill score and $d(s)$ the
position difficulty,
\begin{equation}
\kappa(t,s) \;=\; \max\bigl(-1,\; -\lvert \varsigma(t) - d(s) \rvert\bigr),
\qquad
\kappa(t,s) \geq -0.4 \ \text{ whenever } \varsigma(t) \geq d(s).
\end{equation}
The floor encodes that being overqualified is not a defect.

\begin{definition}[Lexicographic order]
For $\mathbf{u}, \mathbf{v} \in \mathbb{R}^n$, $\mathbf{u} \lex \mathbf{v}$
if and only if there exists $k$ with $u_k > v_k$ and $u_i = v_i$ for all
$i < k$.
\end{definition}

Candidates are compared by $\lex$ on $\score$, with ties broken first by the
member's tie-break rank and then by identifier. Since identifiers are
unique, the comparison is total and the algorithm therefore deterministic.

\begin{remark}[Integrality of $\rho$]
\label{rem:quantize}
The rotation component is integer-valued by construction. This is
deliberate. Under a strict lexicographic order, a real-valued dominant
component makes exact ties a measure-zero event, and the subordinate
components are never consulted: two boards differing by $10^{-4}$ in
rotation would be treated as rotation-distinct, and preference would never
speak. Quantization creates a band of genuine rotation-equivalence within
which the lower objectives decide. The granularity of that band is the one
tuning decision this formulation genuinely requires.
\end{remark}

\begin{remark}[Why not a weighted sum]
\label{rem:bigm}
A scalarization $\Phi = w_1 f_1 + w_2 f_2 + w_3 f_3$ reproduces
lexicographic behaviour only if each weight dominates the entire achievable
range of everything beneath it,
\begin{equation}
\label{eq:bigm}
w_i \;>\; \sum_{j > i} w_j \cdot \bigl(\max f_j - \min f_j\bigr) .
\end{equation}
Condition \eqref{eq:bigm} depends on the \emph{ranges} of the subordinate
objectives. It is therefore invalidated by any retuning that widens a lower
objective, by any unclamped sum, and by the accumulation of several
simultaneously matching preference rows. The failure is silent: no
constraint is violated and no exception is raised; the priority ordering
simply inverts. Lexicographic comparison has no analogue of
\eqref{eq:bigm} to violate, because no arithmetic is performed across
objectives.
\end{remark}

\section{The night algorithm}
\label{sec:night}

Let $G_r = (T \cup S, E_r)$ be the bipartite \emph{eligibility graph} at
relaxation level $r$, with $\{t,s\} \in E_r$ exactly when
$\canplace(t,s,r)$. Coverage is a property of matchings in $G_r$; the
remaining objectives select among them.

\begin{algorithm}[htbp]
\caption{$\nightsolve(\mathcal{I}, \zeta)$}
\label{alg:night}
\begin{algorithmic}[1]
\Require instance $\mathcal{I}$, seed $\zeta$
\Ensure assignment $\mu$, provenance, scorecard, certificates

\Statex \textbf{Phase A --- feasibility census}
\State $E_0 \gets \{\, \{t,s\} : s \in \mathcal{R},\ \canplace(t,s,0) \,\}$
\State $M^{*} \gets \textsc{MaxMatching}(G_0)$
  \Comment{Hopcroft--Karp}
\ForAll{$s \in \mathcal{R}$ unmatched in $M^{*}$}
  \State $\mathrm{cert}(s) \gets \textsc{DeficientSet}(G_0, s)$
    \Comment{Theorem~\ref{thm:hall}}
\EndFor

\Statex \textbf{Phase B --- coverage}
\State $\mu \gets \{\, s \mapsto \iota(s) : s \in \mathrm{lock} \,\} \cup
        \{\, s \mapsto \mu_0(s) : s \in \mathrm{pin} \,\}$
\State \textsc{Augment}$(\mu,\ E_0$ restricted to unassigned members$)$
\If{some $s \in \mathcal{R}$ remains open}
  \State \textsc{Augment}$(\mu,\ E_0)$ along shortest augmenting paths
    \Comment{may displace pinned members}
  \State record each displacement as \textsf{moved-for-coverage}
\EndIf
\If{some $s \in \mathcal{R}$ remains open}
  \State $E_1 \gets E_0 \cup \{\, \{t,s\} : \canplace(t,s,1) \,\}$
  \State \textsc{Augment}$(\mu,\ E_1)$ for open required positions only
  \State tag each $E_1 \setminus E_0$ edge used with
         \textsf{relaxation}$=G_7$
\EndIf

\Statex \textbf{Phase C --- selection}
\ForAll{$s \in \mathcal{R}$ in fill order, $s \notin \mathrm{lock}$}
  \State $C \gets \{\, t \text{ unassigned} : \canplace(t,s,r(s))
          \wedge \textsc{Feasible}(\mu, s \mapsto t) \,\}$
  \State $t^{*} \gets \lexmax_{t \in C} \score(t, s, \mu)$
  \State $\mu(s) \gets t^{*}$;\ \ append $s$ to $W(t^{*})$
\EndFor
\ForAll{$s \in \mathcal{O}$ in fill order}
  \State as above, with $r = 0$ and no displacement permitted
\EndFor

\Statex \textbf{Phase D --- refinement}
\For{$k = 1$ \textbf{to} $K$ (in seeded order)}
  \State propose a swap or relocation $\mu'$ with both endpoints admissible
         at $r = 0$, locks untouched, $\mathrm{cov}(\mu') =
         \mathrm{cov}(\mu)$
  \If{$\Sigma\score(\mu') \lex \Sigma\score(\mu)$}
    \State $\mu \gets \mu'$
  \EndIf
\EndFor

\Statex \textbf{Phase E --- advisory overlay (optional)}
\State obtain proposed overrides from the language model, which is given the
       rules, the per-position candidate rankings, and tools calling the
       \emph{same} $\canplace$ and $\score$
\State accept an override only if $\canplace(\cdot,\cdot,0)$ holds,
       $\mathrm{cov}$ is unchanged, and $\Sigma\score$ does not regress

\Statex \textbf{Phase F --- verification}
\Repeat
  \State $V \gets$ violations of Invariants \textbf{I1}--\textbf{I5}
  \State revert positions implicated in $V$ to their Phase-C values
\Until{$V = \emptyset$}
\State \Return $\mu$, provenance, scorecard (with $V$), certificates
\end{algorithmic}
\end{algorithm}

Two structural points distinguish Algorithm~\ref{alg:night} from a greedy
fill.

First, Phase~B decides \emph{which} positions can be filled before Phase~C
decides \emph{who} fills them. A greedy walk conflates these, and in doing so
commits early choices that can strand a later position: if the only member
eligible for a restroom has already been consumed by a zone, a greedy walk
has no move that recovers. The augmenting path of Phase~B is exactly that
move---it relocates an already-placed member through a chain of
reassignments---and it is found automatically rather than written as a
special case.

Second, Phase~C's oracle $\textsc{Feasible}(\mu, s \mapsto t)$ tests whether
fixing $t$ at $s$ still admits a maximum matching on the remainder. Without
it, an optimal-looking local choice could destroy the coverage guarantee
established in Phase~B.

\section{Correctness}
\label{sec:correctness}

\begin{theorem}[Deficiency certificate]
\label{thm:hall}
Let $G_0$ be the eligibility graph and $\nu(G_0)$ the size of a maximum
matching. Then
\begin{equation}
\lvert \mathcal{R} \rvert - \nu(G_0)
\;=\; \max_{A \subseteq \mathcal{R}}
      \bigl( \lvert A \rvert - \lvert N(A) \rvert \bigr),
\end{equation}
where $N(A) \subseteq T$ is the neighbourhood of $A$. Any maximizing $A$ is a
witness that $\lvert A \rvert - \lvert N(A) \rvert$ positions of $A$ must go
unfilled.
\end{theorem}

\begin{proof}
This is the deficiency form of Hall's theorem, equivalent to
K\H{o}nig's theorem on bipartite graphs. The maximizing set $A$ is
recoverable from the alternating-path structure produced by the matching
algorithm: take the unmatched required positions together with all positions
reachable from them by alternating paths.
\end{proof}

Theorem~\ref{thm:hall} is the algorithm's explanation facility, not merely
its analysis. When five women's restrooms draw on only three eligible
members, the returned set $A$ names those five positions and the deficiency
is two. This is a proof of impossibility localized to the responsible
positions, and it subsumes special-cased headcount arithmetic: a gender
shortfall is not a separate calculation but an instance of deficiency.

\begin{proposition}[Coverage optimality]
\label{prop:cov}
Let $r^{*}$ be the highest relaxation level used by Phase~B. Then
$\mathrm{cov}(\mu)$ is maximum over all assignments satisfying the gates at
levels $\leq r^{*}$ and respecting $\mathrm{lock}$.
\end{proposition}

\begin{proof}
Locked positions contribute forced edges; contract them, which reduces the
problem to matching on the residual graph. By Berge's lemma a matching is
maximum if and only if it admits no augmenting path. Phase~B searches for
augmenting paths over the full residual graph---first among unassigned
members, then, if positions remain open, over paths that displace pinned
members---and terminates only when none exists. Hence the matching is
maximum at level~$0$. If required positions remain open, the edge set is
extended to $E_1$ and the search repeated, yielding a maximum matching at
level~$1$. Maximality at $r^{*}$ follows.
\end{proof}

\begin{proposition}[Selection preserves coverage]
\label{prop:oracle}
If Phase~C admits only candidates passing $\textsc{Feasible}$, then on
termination $\mathrm{cov}(\mu)$ equals the value established in
Proposition~\ref{prop:cov}.
\end{proposition}

\begin{proof}
By induction on the fill order. The hypothesis is that after $k$ selections
the partial assignment extends to a maximum matching. For $k = 0$ this is
Proposition~\ref{prop:cov}. For the inductive step, $\textsc{Feasible}$
admits $t$ at $s$ exactly when $\mu \cup \{s \mapsto t\}$ extends to a
maximum matching; such a $t$ exists because the maximum matching guaranteed
by the hypothesis assigns \emph{some} member to $s$, and that member passes
the test. Selecting any admitted candidate preserves the hypothesis.
\end{proof}

\begin{proposition}[Priority soundness]
\label{prop:lex}
Let $\mu, \mu'$ be assignments with $\mathrm{cov}(\mu) >
\mathrm{cov}(\mu')$. Then $\mu$ is preferred, irrespective of $\rho, \pi,
\kappa$. Analogously, if coverage is equal and $\Sigma\rho(\mu) >
\Sigma\rho(\mu')$, then $\mu$ is preferred irrespective of $\pi, \kappa$.
\end{proposition}

\begin{proof}
Immediate from the definition of $\lex$: comparison is decided at the first
index at which the vectors differ, and no subsequent component is examined.
\end{proof}

Proposition~\ref{prop:lex} is trivial as mathematics and load-bearing as
engineering: it is exactly the guarantee that condition \eqref{eq:bigm}
fails to provide robustly. Under lexicographic comparison the ratified
ordering \eqref{eq:hierarchy} holds by the type of the comparison, and is
therefore invariant under retuning of any component's internal weights.

\begin{proposition}[Determinism]
$\nightsolve$ and $\weeksolve$ are pure functions of
$(\mathcal{I}, \zeta)$.
\end{proposition}

\begin{proof}
Candidate comparison is a total order (lexicographic on $\score$, then on
the tie-break pair, whose second element is unique). All iteration orders
derive from $\zeta$ by a fixed generator. No component reads ambient state:
in particular the rotation component is computed against an explicit
\textit{before-night} cursor rather than the wall clock.
\end{proof}

\section{The week algorithm}
\label{sec:week}

The week solve is the night solve folded over the week, with a shared
rolling history, followed by a cross-night polish. It is not a second
algorithm; in particular it invokes the same gate, so no constraint enforced
on a single night can be evaded across a week.

\begin{algorithm}[htbp]
\caption{$\weeksolve(N, \mathcal{I}, \zeta)$ for nights $N = (n_1, \dots,
n_7)$}
\label{alg:week}
\begin{algorithmic}[1]
\State $W \gets \emptyset$
\ForAll{$n \in N$ in chronological order}
  \State $\mathcal{I}_n \gets$ instance for $n$ with history
         $\mathrm{merge}(H, W)$
  \State $\mu_n \gets \nightsolve(\mathcal{I}_n, \zeta \oplus n)$
  \State $W \gets W \cup \mu_n$
    \Comment{planned nights are visible to later nights}
\EndFor
\For{$k = 1$ \textbf{to} $K_W$ (in seeded order)}
  \State propose exchanging placements across nights $n_i \neq n_j$, both
         endpoints admissible at $r = 0$, coverage of both preserved
  \If{$\textsc{WeekScore}$ improves lexicographically}
    \State accept, and \textbf{rebuild} the context of every night after
           $\min(n_i, n_j)$
  \EndIf
\EndFor
\State \Return $(\mu_n)_{n \in N}$, week scorecard, relaxation ledger
\end{algorithmic}
\end{algorithm}

\noindent
The week score governing the polish loop of Algorithm~\ref{alg:week} is
\begin{equation}
\textsc{WeekScore} = \Bigl(\textstyle\sum_n \mathrm{cov}(\mu_n),\;
-\!\sum_n \mathrm{rep}(\mu_n),\;
\sum_n \Sigma\rho,\; \sum_n \Sigma\pi,\; \sum_n \Sigma\kappa \Bigr),
\end{equation}
with $\mathrm{rep}$ counting within-week area repeats. The context rebuild
that accompanies each accepted exchange is required for soundness: an
exchange alters the history that later nights were evaluated against, and
evaluating a night against a board that no longer exists admits exchanges
whose true downstream effect is negative.

\section{The placement matrix as a derived view}
\label{sec:matrix}

The placement matrix records, per member and per zone, exposure counts over
trailing calendar windows. We observe that it is a pure function of history
and the clock:
\begin{equation}
\label{eq:matrix}
\mathcal{M}(t, a) \;=\; \bigl( c_{28}(t,a),\; c_{56}(t,a),\; \ell(t,a)
\bigr),
\qquad
c_{w}(t,a) \;=\;
\bigl\lvert \{\, e \in H(t) : \alpha(e) = a,\ \delta(e) \leq w \,\}
\bigr\rvert,
\end{equation}
where $\delta(e)$ is the age of event $e$ in days and $\ell(t,a)$ the most
recent night on which $t$ occupied $a$. Consequently
\begin{equation}
\mathcal{M} \;\equiv\; f(H, \mathrm{now}),
\end{equation}
and materializing $\mathcal{M}$ is a caching decision rather than a design
one. At the operative scale---order $10^3$ history rows---evaluating
\eqref{eq:matrix} on read is submillisecond, and the materialized form
should be regarded as an optimization to be justified, not a default.

The distinction matters because a materialized $\mathcal{M}$ admits states
that \eqref{eq:matrix} cannot express. A stored count can disagree with
history, either because a mutation path failed to trigger a rebuild, or
because trailing windows were computed at a past instant and no longer
correspond to the present one. Neither failure is representable when
$\mathcal{M}$ is evaluated on read. If materialization is retained, the
governing invariant is that $\mathcal{M}$ be reconstructible from $H$ alone,
identically, at any time; a rebuild triggered by every mutation that touches
history, together with a periodic sweep, then bounds staleness while
eliminating incorrectness.

\section{An equivalent integer program}
\label{sec:ilp}

The same problem admits a compact integer formulation, useful as an
independent oracle against which the staged algorithm can be validated.

\begin{align}
\text{variables}\quad
& x_{ts} \in \{0,1\} && \text{$t$ occupies $s$} \\
& y_s \in \{0,1\} && \text{$s \in \mathcal{R}$ left open} \\
& z_{ts} \in \{0,1\} && \text{edge used under relaxation } G_7 \\[0.4em]
\text{subject to}\quad
& \textstyle\sum_{s} x_{ts} \leq 1 && \forall t \in T \\
& \textstyle\sum_{t} x_{ts} + y_s = 1 && \forall s \in \mathcal{R} \\
& \textstyle\sum_{t} x_{ts} \leq 1 && \forall s \in \mathcal{O} \\
& x_{ts} = 0 && \forall (t,s) \text{ failing a Level-0 gate} \\
& x_{ts} \leq z_{ts} && \forall (t,s) \text{ failing only } G_7 \\
& x_{\iota(s)s} = 1 && \forall s \in \mathrm{lock}
\end{align}

The objective is optimized in stages, each stage constraining the next:
\begin{enumerate}[itemsep=0.1em,topsep=0.25em]
  \item minimize $\sum_{s \in \mathcal{R}_1} y_s$, then
        $\sum_{s \in \mathcal{R}_2} y_s$ \hfill (coverage, by tier)
  \item minimize $\sum_{t,s} z_{ts}$ \hfill (fewest relaxations)
  \item maximize $\sum_{t,s} \rho(t,s)\, x_{ts}$ \hfill (rotation)
  \item maximize $\sum_{t,s} \pi(t,s)\, x_{ts}$ \hfill (preference)
  \item maximize $\sum_{t,s} \kappa(t,s)\, x_{ts}$ \hfill (skill)
\end{enumerate}

Staged optimization---each stage fixing the optimum of its predecessor as a
constraint---realizes \eqref{eq:hierarchy} without weight arithmetic, and is
the standard formulation of preemptive goal programming. Since $\pi$ depends
on the assignment through pair affinity, that term is either linearized with
auxiliary products or, in practice, evaluated in the final stage over a
restricted candidate set.

\begin{remark}
At $\lvert T \rvert \approx 25$ and $\lvert S \rvert \approx 34$ the
program has fewer than a thousand binaries and solves in milliseconds.
Running it as a shadow validator alongside the staged algorithm yields a
disagreement certificate---a concrete board on which the two differ---which
is a considerably stronger signal than an assertion failure.
\end{remark}

\section{Invariants}
\label{sec:invariants}

\begin{description}[itemsep=0.15em,topsep=0.3em,leftmargin=3.4em,
                    font=\normalfont\bfseries]
  \item[I1] $\mu$ is injective: no member holds two positions.
  \item[I2] Locked positions are identical across any run.
  \item[I3] No output placement violates a Level-0 gate.
  \item[I4] $\mathrm{cov}(\mu)$ attains the matching bound of
            Proposition~\ref{prop:cov}: coverage is optimal, not
            best-effort.
  \item[I5] Every $G_7$ waiver appears in the relaxation ledger, and the
            number of waivers is minimized.
  \item[I6] $\nightsolve$ and $\weeksolve$ are pure in
            $(\mathcal{I}, \zeta)$.
  \item[I7] $\mathcal{M} \equiv f(H, \mathrm{now})$, over a single history
            store with singular $(\textit{night}, \textit{position})$
            ownership closed under every mutation path.
  \item[I8] $\canon$ is total over production key forms; unrecognized keys
            fail closed.
  \item[I9] Explanations---score vectors, deficiency certificates,
            relaxation ledgers---are byproducts of solving, not
            reconstructions after it.
\end{description}

Invariants \textbf{I1}--\textbf{I5} are checked in Phase~F and re-checked
after any repair, so a result reporting no violations has been shown to have
none. \textbf{I3} and \textbf{I8} are amenable to property-based testing over
randomized rosters, configurations, and seeds.

\section{Complexity}

Let $n = \lvert T \rvert$, $m = \lvert S \rvert$, $E = \lvert E_0 \rvert
\leq nm$.

\begin{center}
\begin{tabular}{@{}llll@{}}
\toprule
Phase & Operation & Complexity & At $n{=}25, m{=}34$ \\
\midrule
A & Hopcroft--Karp & $O(E\sqrt{n+m})$ & ${\sim}10^{4}$ ops \\
A & Deficiency certificates & $O(E)$ per open position & negligible \\
B & Augmentation & $O(E\sqrt{n+m})$ & ${\sim}10^{4}$ ops \\
C & Selection with oracle & $O(m \cdot n \cdot E)$ worst case &
    ${\sim}10^{6}$ ops \\
D & Refinement & $O(K)$, budget-bounded & tunable \\
F & Verification & $O(m)$ per pass & negligible \\
\bottomrule
\end{tabular}
\end{center}

The whole solve completes well within ten milliseconds. No approximation is
required at this scale: provable coverage optimality and full determinism
are available at no practical cost, and any heuristic in the implementation
should be justified by explanatory or incremental-update requirements rather
than by performance.

\section{Relation to the deployed implementation}

The deployed system is staged in the same shape---plan, optimize, advise,
verify---and that shape is worth preserving, since it is the source of the
provenance that makes engine output auditable by the operator. The
formulation above substitutes for its internals as follows.

\begin{center}
\begin{tabular}{@{}L{0.29\textwidth}L{0.30\textwidth}L{0.30\textwidth}@{}}
\toprule
\textbf{Deployed} & \textbf{Specified here} & \textbf{Consequence} \\
\midrule
Greedy fill with rescue and backfill passes &
Maximum matching, then selection under a feasibility oracle &
Coverage provably optimal; stranding unrepresentable
(Prop.~\ref{prop:cov}, \ref{prop:oracle}) \\[0.3em]
Scalar weighted sum with hard gates &
Lexicographic vector $(\rho, \pi, \kappa)$ &
Ordering \eqref{eq:hierarchy} holds by construction
(Prop.~\ref{prop:lex}, Rem.~\ref{rem:bigm}) \\[0.3em]
Three distinct relaxation behaviours &
One two-rung ladder inside $\canplace$ &
Single semantics; every waiver ledgered \\[0.3em]
Multiple history stores and key vocabularies &
One store, $\canon$ at the boundary &
Rotation false-negatives eliminated as a class
(Def.~\ref{def:boundary}) \\[0.3em]
Materialized matrix with refresh choreography &
Derived view; materialization optional &
Staleness bounded, incorrectness unrepresentable
(\S\ref{sec:matrix}) \\[0.3em]
Repair without re-verification &
Verify--repair--re-verify to fixpoint &
A clean report is a proof, not an assumption \\[0.3em]
Feasibility by headcount arithmetic &
Deficiency certificates &
Exact, per-position, gender split subsumed
(Thm.~\ref{thm:hall}) \\
\bottomrule
\end{tabular}
\end{center}

\vspace{1em}
\noindent
Two companion documents complete the set:

\begin{description}[itemsep=0.2em,topsep=0.3em,leftmargin=1.2em]
  \item[\normalfont\texttt{PLACEMENT\_ARCHITECTURE.md}] --- a narrative
        description of the system as built, including the persistence layer
        and the explanation surfaces.
  \item[\normalfont\texttt{PLACEMENT\_REVIEW\_2026-07-18.md}] --- a
        catalogue of the specific divergences between the deployed system
        and this specification, with severity ranking and a suggested
        remediation order.
\end{description}

\end{document}
